% ================================================================
%  18.5.4 注 (c1) 之尾 — (d) — (e) — (f) — c3 — c4 — Epilogue outre-tombe
%  对应原书 p.813–836（PDF 物理页与印刷页同）
%  源文: _work/chunks/chunk01–07.md
% ================================================================

% ---------------- p.813 ----------------
作为指向同一方向的第三个迹象:我曾听说,克拉斯纳(五六十年代的巴黎
数学圈里人尽皆知的一位怪才:他家里收容着一支猫大军,裹着一件厚重的俄式大衣,永远一副
乐呵呵的表情,在各讨论班之间走来走去……)——就是这样一位克拉斯纳,据说在非阿基米德
赋值域上“做解析延拓”。我所知仅此而已;我甚至不能确定可曾遇见过哪一位读过克拉斯纳这
方面工作的人——但这件事本身已够引人好奇。应当说,“解析延拓”这个词本身可没有让我心
跳加快的魔力(毋宁说相反,它叫我想起学生时代一些索然无味的记忆……);可一旦隐约察觉
到需要一类新的几何对象,它就难免让人心头一震……

回到Remmert——既然他的记性糟糕到这般地步,Tate的原文本(他还自夸握有一份)总该能替
他唤回记忆。在笔记里,Tate对我在这套理论的构想中\dmark\origpage{1066}所扮演的角色没有丝毫讳言\fnote{795}{(*)}{自Tate与我、他的家人与我的家人亲密往来的那些遥远日子以来,二十多年已经过去。多少年了,再没有他的只字音讯。这本书印出来,我的前学生与旧友当中没有人会看不到这篇引言里对我本人的这番抹除——可我也从未听说他们中哪一个为此动过声色。此一时,彼一时……},其中写道(凭记忆
引述):他“完全忠实地”遵循一位“总施工”(maître d'oeuvre——这是就通过“拼
合碎片”来构造该概念的方法而言),而这位总施工正是他从我这里得到的。此外我还为他提供
了某种“砌块”(或者说:在受限形式幂级数代数中进行“局部化”的“方法”),以应形式概形
纤维之需。他又以第二类“碎片”(或“方法”)补足这最初的一批,两者在某种意义上互为补充。

若不是我手里有“广义空间”这条主线(它后来成为\textbf{topos}),这个新概念恐怕根本不会问世
——平展上同调不会,晶体上同调不会,尾随其后涌现的许多东西也不会,甚至连那道最时兴的
“奶油甜点”、大名鼎鼎的 D-模也在内;这条主线的理论当时尚待建立,但四年来已被人隐隐预感。
正是这条直觉,为我指明了通往一类“簇”的道路——这类“簇”恰恰\textbf{越出了}寻常拓扑空间
(局部赋环空间)的框架。

自从John~Tate启动\textbf{刚性解析空间的局部理论}那一刻起,同样是我提出并推广
了关于这些新簇的首批关键“整体”定理的表述;这些表述早在第一批基础工作完成之前,就已
经在我心中成形:关于刚性解析空间上相对真概形的代数-解析比较定理;关于刚性解析空间的
真态射 $f$ 的 $R^i f_*$ 的有限性定理——这些问题由Kiehl在随后几年里解决\fnote{796}{(**)}{这里要指出,早在Tate打下刚性解析空间理论最初的地基时,我就清楚:他所处的那个框架仍是暂定的,丝毫不足以穷尽我试图用“刚性解析空间”这个名字表达的直觉内容——正如域上有限型概形穷尽不了“概形”一词所系的直觉。有一条线索指向Tate框架的实质扩充(凡愿意听的,我都讲过……),而提供它的是Tate自己:他在某个拓扑环上写出了一条“Tate泛椭圆曲线”(这个环是形式幂级数环 $\mathbb{Z}[[t]]$ 的一个子环,由那些在复平面开单位圆盘内收敛的级数组成——如果我记忆无误),而这个环显然应当被看作某一个刚性解析空间的“仿射坐标环”,其类型并不在Tate开出的清单之内。自我离去之后,一切基础问题都跌入人人不屑的境地;1962年Tate搭起的这台概念机器从那以后纹丝未动,也就不足为怪了。}\fnote{797}{(***)}{(6月4日)我还曾最先力主:对刚性解析空间,必须引入比Tate所设想的更一般的“点”(Tate的点只取基域的\textbf{有限}扩张为值)。这一必要性,既受代数几何之类比的提示,也出于一个愿望:为与所论刚性解析空间相伴的 topos 的“点”找到一个具体的解释。}。不过
话说回来,照眼下盛行的风气,\dmark\origpage{1067}预想新概念、物色总施工、提出供“真正的数学家”去解决
的问题,都被当作无关紧要之事——往极端里说,简直就是唬人的花招……

% ---------------- p.814 ----------------
不管怎样,这篇引言里没有一字提到我的名字与刚性解析空间有什么相干。克拉斯纳的名字同
样不见踪影——非但如此,Tate的理论还被说成引入了“一种足够丰富的结构,使不可能成为可
能:全不连通域上的解析延拓”——可早在1962年,这项(“不可能”的)解析延拓,十年也罢,
二十年三十年也罢(我说不准),早已是克拉斯纳“注册在案的正业”(若可以这样说)了。汗
牛充栋的参考文献里,克拉斯纳和我都无迹可寻。不过快到引言末尾,我的名字倒是顺带闪现了
一回,在“Grothendieck 拓扑”这个名字里;这个概念所引的是Artin1962年的笔记,
而对 SGA~4 里那份精益求精的打磨功夫(自1963年起,贯穿整个六十年代——却顶着一件显然
不受欢迎的父名……)则傲然无视,完全步了我那班前学生的后尘……;同样只字未提的,是当
年(1966年)我曾为刚性解析空间在晶体上同调的发展中指派的那个角色——那时Remmert(连同
他那些显赫的复分析家同事们)还全无迹象要理会这类所谓的“刚性解析”的怪“簇”(这名字
您给评评理……),它们是几位代数几何学家在自己角落里捣鼓出来的——仿佛复解析空间还
不够那些严肃的分析家与几何学家消遣似的……

\pmark\origpage{1068}关于书中所述这套理论的诞生史,只要是从第一手知情,就能看出这篇引言里招摇着同一种犬
儒做派——它同样流露在那份匿名审稿人对一位无名申诉者的答复里(并蒙这位Remmert先生祝福):显然,在作者们的心目中,这不过又是一桩“礼数”问题,说到底是一种“客
气”,施与不施全凭他们高兴:在新理论的孕育过程中起过关键作用的某某某之名,要不要写进
他们的“历史概述”(原文如此),悉听尊便。在他们看来(恐怕在数学界几乎所有建制中人看
来亦然——他们对此类歪曲照单全收,连眉头都不皱一下……),“历史”并不是\textbf{实际发生过
的事},而是一件可以由僭取书写之权者\textbf{乾纲独断}的东西,或者由一小撮人合议而定:什么有资格
存在,什么有资格算作发生过,全由他们说啦。

这班人就爱对苏联发生过、还在发生的事情津津乐道,而且一桩也不会放过(我知道自己在说
什么),忙着签署什么“捍卫自由”(思想自由之类……)的宣言——那是在\textbf{别人}家里;而在
\textbf{他们自己}掌权的地方,同样的谎言专制一天也没停过。

(6月3日)几天前在前几页里提到克拉斯纳那副独特而可亲的形象时,我心头冒出一个问题:
他还在世吗?他比我年长一两辈,我已有很久(少说十五年,也许二十年)没听人提起他的名字。
人物我还记得清清楚楚,他的名字却要过几秒钟才浮上来。(说实话,这一类事如今常在我身上
发生,年纪不饶人……)克拉斯纳以好客闻名,他的俄国血统本是另一个本可让我们接近的共同
点。可我那时一头扎在数学里,腾不出闲工夫只为“有趣”去交朋友。我们俩对待数学的路数,
想必正处在对跖点上。我们大概在Bourbaki讨论班的间隙闲聊过一两回,数学是肯定没有谈的。
那时真正让我着迷的,除了数学几乎没有别的……

% ---------------- p.815 ----------------
\pmark\origpage{1069}话虽如此,今天我收到Deligne的一张便条,不过寥寥数行,谈一桩无关紧要的实际琐事,
或许是想让自己还留在我的好记忆里(我们之间没有书信往来,算来已有几个月了);又或许是
想找个地方安插一条附言——我冒昧把它抄在这里(想必他会同意):

\begin{quote}
“又及:惊闻克拉斯纳已于半月前去世,不胜哀恸。我一直记得二十来年前他在布鲁塞尔所
做的一场滔滔不绝的长篇报告,我自然是云里雾里,却硬是撑到最后,成了所剩无几的听众之
一。令我吃惊的是,你那张五十年代人物表\fnote{798}{(*)}{这里对我在《收获与播种》第一部分(“Fatuité et Renouvellement”)中的说法有一个明显的误会。那一部分的主旨从来不是勾画一幅数学“五十年代群像”——哪怕只限巴黎圈子,或者Bourbaki周围那一圈。我的根本意图,是去发现我作为数学家的过去。正因如此,我才谈到我与这样那样同事或学生的关系——每当这层关系在我的生命中显得重要,或者能够照见我自己的时候。}里竟没有他——他在那个年代做了不少漂亮事,虽然与Bourbaki的精神格格不入,
而且有一种专把定义写得潦草马虎的天才。”
\end{quote}

这样一来,这里又多出一篇悼词,这一次是为我的一位同穴之人而作。这一篇里,我自信看出
一丝同情,或者毋宁说是某种情感的余光,那种情感往日曾经鲜活。可是,同在我的悼词里一样,
我的朋友皮埃尔决不会松一松牙关,说出——这一次是为了纪念一位一去不返的逝者——他津津
乐道却从不点名的那些“漂亮事”究竟\textbf{是什么}。然而他和我一样清楚:这些“事”为一个
如今正开得繁盛的理论的问世铺了路;而且,出于他或许心知肚明的原因,新主子们乐于抢先把
这位和善、潦草、“马虎”的先行者提前下葬(就葬在我的旁边)——他刚刚离世。他确凿无疑
是一位在超度量域上“做解析延拓”的人物——那时Tate、Remmert或者我还在“做”三角形全等
的种种情形和毕达哥拉斯定理,而我们的朋友皮埃尔还得让妈妈替他擤鼻涕(顺便擦洗……)呢!

% ---------------- (d) La Répétition Générale (avant Apothéose) — p.815–818 ----------------
\subsubsection*{(d) 总彩排(封神之前)}
\tocsubsub{(d) 总彩排(封神之前)}

(4月16日)可我必须回到我身后学生Zoghman~Mebkhout那一连串够味儿的“倒霉遭遇”上来。
\dmark\origpage{1070}1979年6月,当一个籍籍无名的陌生人——口口声声以Grothendieck的思想为宗——亲口告诉他,
那个让Deligne十年前整整一年绞尽脑汁、终未得到满意解答的关键问题\fnote{799}{(*)}{(5月25日)也可能,Deligne早就失去了对“关键”二字的分寸。关于这一点,见注“……以及羁绊”(n\textsuperscript{o}~171(viii)$'$)。}有了一个漂亮的解法时,Deligne的脑子里究竟发生了什么,我毫无概
念。以他一贯的心性,可想而知,他绝不会去祝贺这个年轻人在他Deligne栽过跟头的地方得了
手。可我深深感到,他那掘墓人的心性使他的嗅觉(我曾领教过的惊人嗅觉)失灵到这般地步,
以至于直到今天(六年之后),连他自己也没能把握住那个无名之辈的思想与眼光的真正分量。
他和所有人一样,最终只看见了那道“奶油甜点”:那件人人翘首以待的意外工具,那把砸开
“出了名的难题”的铁钎。可他也曾有一天,把别人传递给他的一整套广阔眼光据为己有——为
的是把那眼光连同生下它的人一并埋葬,再顺手抢走又一件工具,把它也改造成一柄“铁钎”……

我所知的Deligne对Mebkhout定理的最初一次反应,是一封未注明日期的短信(手写),写给Mebkhout,收到于1980年10月10日\fnote{800}{(**)}{这就是那份“以保密为条件提供、我在此不再多置一词……”的文件,“受害者”一注(打字稿第309页)里提到过它。事隔一年,Zoghman慨然允准我在此照录。}:

\begin{quote}
亲爱的Mebkhout:

我已把我那本你的论文寄给Bernstein和Beilinson:他们证明 Kazhdan--Lusztig 猜想用得
着你的结果(我有一份他们工作的俄文摘要,你要的话我可以转寄)。能不能再给我寄
一本?

谢谢。

P. Deligne
\end{quote}

从这封信我推想,Deligne必定已把“好上帝定理”的事告知了那两位苏联数学家,或许还暗示它
可以用来证明那个猜想;要么这是他自己悟出来的,要么当时已有风声,说Brylinski在这个
题目上有想法。Mebkhout那场“点着了”Brylinski的报告,早在1980年1月。Brylinski-Kashiwara
一边,Beilinson--Bernstein一边,两篇文章——用那个没有名字的定理(它的作者更没有名字)证
明了那条著名猜想——\dmark\origpage{1071}收到日期一为1980年12月19日,一为1980年12月8日,前后只差十一天。
巧合?

我甚至冒出过一个念头:Deligne1979年6月就比谁都先知道这件新工具(反正谁也不曾费神去
读那个无名之辈的厚砖头,Deligne也没有)——他为什么不自己想到把它用在那个猜想上,收
获一顶新桂冠,反倒去帮苏联同行摘桂冠?论头脑的敏捷,他不见得输给Brylinski吧?说不定
从那一刻起,他就已经瞥见一种可能:从旁门把“好上帝定理”本身的发明权捞回来——这发明权
(他心里必是这么感觉的)十年前就该是他的;说不定在他看来,正是由于一场不可饶恕的发牌
错误,那个不自量、欠打磨的狂妄后生才僭取了权利,去证明他Deligne早已久久钻研而终无决定
性成果的东西。他最终离成功只差那么一丝丝;换了别人在他流汗无功的地方收割,这可不公道
……但若想收回那件归根结底“理所应当”归于他的东西(依照某个人上人圈子里终于通行起来
的不成文法——他自认是那个圈子的中心和大佬……),那他就得以全然不同的细腻手法来操盘,
不能妄图一口吞下太多\fnote{801}{(*)}{这当然只是揣测:关于那条著名“对应”的据为己有之念,是否早在Deligne初闻其事时就已存在。就我自己而言,我深信不疑。诚然,上引那封信看上去会给人相反的揣测。而在我看来,这仍是那种挑衅姿态的又一个信号:他Deligne根本不必有半点顾忌——反正对方不过是个无名之辈,孤身一人对全世界,\textbf{横竖动弹不得};而他Deligne尽可以“自甘出丑”,正如他尽可以借“反常层”(\emph{faisceaux pervers})这个挑衅性的命名,以一种象征性却刺眼的方式,宣告自己真实的心性。关于这一点,见注“反常”(n\textsuperscript{o}~76),以及(在心理情境十分相近、却不那么极端的场合)注“玩笑——或‘权复形’”(n\textsuperscript{o}~83)。}。

话虽如此,Zoghman吃过与Kashiwara、与Brylinski打交道的怪插曲的亏,已有戒心,他审慎行事,决定亲自向Beilinson、Bernstein两位先生通报那个定理——Deligne说他们用得着它——以防Deligne这么大的
人物在谈起定理时,忘了提一句谁是它谦卑的作者。时机正好:下个月,1980年11月24日至28
日,莫斯科有一场“广义函数及其在数学物理中的应用会议”。Mebkhout在会上报告了他的定理,
已刊行的\dmark\origpage{1072}题目是“高维的Riemann--Hilbert问题”;他特意与Beilinson、Bernstein当面交谈,把自己
结果的来龙去脉原原本本讲给他们听。

他来的正是时候。会后不过十天,两位作者就寄出了他们关于 Kazhdan--Lusztig 的工作,以《法
国科学院院报》快报的形式(第292卷,1981年1月5日,第I辑--15):“群论——$g$-模的局部
化”,Alexander~Beilinson与Joseph~Bernstein署名,\textbf{由Pierre~Deligne转交}。理所当然,他们的
手稿上只字未提Mebkhout之名——显然,Deligne把那个无名之辈忘了个一干二净,尽管他明明把
此人的论文转寄给了他们,其目的正是……?谁懂,谁就懂吧!Mebkhout费了九牛二虎之力,才
让Beilinson(“两人中更老实的那一个”,Mebkhout一本正经地向我担保)相信:他们书目里引的
Kashiwara--Kawai那篇文章里,唯独没有那件“构造”(此处替代那句说了八百遍的“对应”)——而这件
“构造”,他们同所有人一样,也只敢用暗语提及。(多半是Deligne,一边把那篇无名之辈的论文
转寄给他们——想要的结果清清楚楚就在里面\fnote{802}{(*)}{(4月17日)论文里至少有一个非常接近的结果,尽管Beilinson--Bernstein(以及Brylinski-Kashiwara)所用的那套表述在其中并未逐字出现。细节见同日一条脚注(打字稿第1047页注(**))。}——一边必定向他们递了话:若非要给个参考文献,引Kashiwara的某一篇文章或许更稳妥些,哪一篇倒在
其次,反正没人会去查证到底。)对于那位亲自到场的无名之人,他们倒也当面许诺:会想着这
事的,会在Kashiwara那一处把引用改正过来。抱歉抱歉——我朋友Zoghman的倒霉史,确实老是那一套!

在这几位杰出作者的呈文里——由Deligne转交(他的信我刚才才照录,写于仅仅一个月之前)
——\textbf{Mebkhout的名字只字不见}。顺便说,Kashiwara的名字也不见(我已经瞥见那儿支起了一只耳朵
……)。不过,在呈文的最后一部分(证明 Kazhdan--Lusztig 的那一段)里,倒有一处偷带的两重
引注,引一件“\textbf{见于}[4]、[5]\textbf{的构造}”\fnote{803}{(**)}{读者不妨细品“\textbf{见于}……的构造”这一说法的含糊其辞:这件“构造”(或“对应”、或“关系”……)究竟出自谁手,问题被完全悬置;而仅仅六个月之后,在那场著名的讨论会上(见注“变戏法者”,n\textsuperscript{o}~75$'$),这个问题将以人们熟知的精湛手法得到“解决”:从Beilinson--Bernstein-Deligne的文章里,读者将获悉,那电报式的引注[4][5](在两处地方——那两处,想必那件构造恰好(托天之福)“有所陈述”)纯属客套,而“对应”的光辉父亲,果然是大家猜想的那一位……不过,撇开我刚提到的这套魔术手法不谈,引用一个新定理——又深又难——却只写成“见于……的构造”,仿佛那不过是件寻常“构造”,碰巧被丢在那里,作者们又碰巧选中它,拿来支撑他们辉煌的证明——这本身就已是一桩骗局。我在这里认出的,正是“SGA~4$\frac{1}{2}$--SGA~5”那场手术的同一副心肠:那场手术的要领,是(顺带)提一句“平展上同调形式‘有所陈述’于 SGA~4 与 SGA~5”(以及那位光辉作者不得不从中剔除的“废话矿渣”),随后挽起袖子,做作地开始“干\textbf{正经}数学……”。(5月25日)关于这种“新风格”,另见注“恭贺——或新风格”(n\textsuperscript{o}~169$_9$)。},\dmark\origpage{1073}这件“构造”(您准猜着了!)不是别
的,正是那位当值无名之辈那个从未被命名、其人更不曾被叫出名字的函子。参考文献[4]是Kashiwara的
一篇文章(临时的替身爸爸)。那篇文章里(Kawai-Kashiwara那篇也一样——它被扔进了盈亏账)自
然没有任何东西,哪怕远一点儿,像这几位作者所说的“构造”;况且那篇文章是1975年的\fnote{804}{(*)}{核实过了:就是前文引过的Kashiwara那篇文章,他在其中证明了他的可构造性定理——这个定理当然对定义那些“好上帝函子”起着关键作用(不过在1980年那阵抢购潮之前,除了Mebkhout,谁也不曾梦想过这些函子)。佯装把Kashiwara这个定理(无人打算跟他争)与“好上帝定理”混为一谈——后者无比深刻,层次完全不同——这是一桩粗劣的骗局。就证明而论,“好上帝定理”动用了Heisuke~Hironaka式奇点消解的全部威力。就“哲学”而论——这还要重要得多——它在同调形式主义内部架起了先前缺失的桥梁,沟通拓扑、代数与分析(算术尚待来日——前提是,我眼看着正在掘墓的某些先生,总有一天能找回他们健全的心智……)。},比莱苏什讨论会上某位无名之辈的报告早了差不多五年——正是那场报告让这位
Kashiwara先生开了窍:把“导出范畴”这个词说出口,毕竟不算傻事,而且可以依仗“强者即有理”这条再简单不过的法
则,把别人流汗挣来的功劳据为己有。至于参考文献[5],那是Mebkhout1979年9月在莱苏什讨
论会上的报告——恰恰就是在那场报告上,Kashiwara先生学到了:导出范畴原来是有用的——而且
除了用来坑害一位被他的大佬和前辈们晾在一边的无名之辈,还能派上别的用场……

\pmark\origpage{1074}同Brylinski-Kashiwara那篇文章一样,这里也没有任何东西,能让一个真正“不入局”的读
者起半点疑心:若不是那件新兴而天赐的工具现世——它被“见于……的构造”这句托词悄悄
藏起——这篇光辉的呈文根本不会问世。我还认出了那套屡试不爽的“溺毙鱼儿”术\fnote{805}{(*)}{关于这种号称“同化稀释法”的招数的其他例子,见子注“\textbf{正经}数学……”(n\textsuperscript{o}~169$_5$),即打字稿第885页脚注(*)。},号称“稀释法”:把要抹掉的那个人(却又想
给自己留“叫停”的余地,好在必要时说自己引过他……)\textbf{配对}到另一个与此事毫不相干、
或作用微不足道的人身上——仿佛在此处(在字里行间,却又明白无误地)说:我们把那个无名
小卒放在这儿(纯粹出于客套,也架不住他一再坚持),他和那件大名鼎鼎的“构造”(按新近确
立的共识,只能用暗语谈它,把它当作人人皆知的东西……)的干系,就跟1975年那篇文章的干
系一样多——那一年,“上流世界”还没有谁肯屈尊说出“导出范畴”这个词(顶多当笑话讲讲
……)。

\subsubsection*{(e) 冤大头买卖——或木偶戏}
\tocsubsub{(e) 冤大头买卖——或木偶戏}

我不后悔花这番功夫——既为我自己的备忘,也为哪位对这桩事感兴趣的数学读者——在这里
逐一清点围绕当值无名之辈那条定理的三桩前奏骗局。这三桩骗局,一桩是Kashiwara家的土产,一桩是
Brylinski-Kashiwara的土产(还劳了一位始终匿名的审稿人从旁襄助),一桩是Beilinson--Bernstein
的土产——而Deligne,则在幕后\fnote{806}{(**)}{(6月5日)Deligne“在幕后”扮演的角色,至少在第三段插曲里是明摆着的;第二段插曲,也有种种强烈的迹象指向同一方向。但看起来,Kashiwara早在1978年就已经在为他自己“打响第一枪”(就围绕Mebkhout著作的那些骗局而言)——当时(看来)Deligne还什么都蒙在鼓里。关于这一点,见本注 c) 部分(“入门费——或有为青年”),即打字稿第1050页脚注(*)。}。三桩在风
格上如出一辙,齐整得惊人,无须我再多着墨。这种风格,在我对\emph{葬礼}\fnote{807}{(***)}{关于这种风格,特别可看前已引过的注“恭贺——或新风格”(n\textsuperscript{o}~169$_9$)的结尾。}所
作的漫长调查一路之上,我早已领教到倒尽胃口;而在我那位天资最出众的学生——正是那位Pierre~Deligne——1968年的一篇文章里,它已经以刺眼的方式预先现出了形\fnote{808}{(****)}{见注“驱逐”(n\textsuperscript{o}~63)的开头。}。而这层情形也足以\dmark\origpage{1075}提醒我的记性:当年面对Deligne以及另一些在我看来天资出众的人,我抱着一种模棱两可、曲意逢迎的态度,
对今日铺展在我眼前的这股败坏,我也不是没有添上我自己的一份。

现在也清楚了:1981年6月那场“反常层会议”(Colloque~Pervers)的\emph{封神},离我们方
才清点的第三段插曲不过半年,绝非从天而降。说也奇怪,据我所知,这场讨论会是我离去之后头
一场、也是唯一一场专门(嘴上自然绝口不提,实际上却毫不含糊)用来把“Grothendieck数学”的
某个侧面重新发掘出土的讨论会——起因是一件猝然降临的新工具,谁料它竟无可替代。这件工
具,唯有在一种看待事物的路数里才活得下去——而时尚的种种“共识”,老早以前就把这种路数
归档为过时,还搭上几分隐隐的滑稽\fnote{809}{(*)}{至于这些“时尚共识”背后起作用的心理机制——它们与面对某种数学研究路数时某种“出自本能的”排斥反应彼此叠合——见前已引注“天赐的机缘——或封神”与“拒认——或召回”(n\textsuperscript{o}~151、152)。}。而仗着我那位天资夺
目的前学生的特殊天才,事情起了一场古怪的回旋:这条被他本人、被所有人弃绝的路数,竟在需
求的推动之下于事实中得到辉煌的证实;而恰恰又是借着同一场讨论会,这场证实成了机缘,把那
位已故而未被提名的大师——连同那位身后的学生(同样无人提名,他曾有幸(或不幸……)让这一大帮上流人物为他挪步前来)——彻底而最终地下葬。

% ---------------- p.819 ----------------
这场讨论会,可不是从天上掉下来的,决不是。我的朋友Pierre~Deligne有个特点(别的且不
论):他懂得等待,也懂得抓住有利的时机。围绕那道“奶油甜点”的三段插曲,加上对当值无名
之辈的提及早已被删削得几乎一干二净,分明向他表明:时机已经熟透,该神不知鬼不觉地、带着他那副
招牌式的笑容可掬与温文尔雅,把那件横竖都理应归他名下的东西收进兜里了。我猜他跟Verdier做过一番周密的串通:得让Verdier明白,如今时候已到,该锣鼓喧天把导出范畴和那份早已被
拒认的“父亲名分”从土里重新刨出来;与此同时,又把那个模糊的无名之辈、连同那位谢世多年的大
师,一并埋进舞台脚灯的光焰底下(就怕哪位不识相,偏偏记起:那些俨然“最新时尚”的美好
玩意儿,那位大师当年原也出过一份大力……)。

\pmark\origpage{1076}让Kashiwara给某条从未被点名的“好上帝定理”当个偷偷摸摸的爹,这一套一度还对付得过去
;而且恰恰也只在这样一个期间对付得过去:大家讲好了,那个定理既不会被点
名,也不会见诸文字。说到底,Kashiwara本人对这条定理也未必真有多么上心——这条定理,说不定他懂
的比Verdier本人还少;八成是机会凑手、积习相帮,他顺手就把它
揣进了腰包,仿佛出于无心。至于Deligne,那位懂得等待的人,他清楚得很:这条定理不会永远是
一条既无地址、也无姓名的定理。总而言之一句话:这是一条\textbf{在给自己寻找父亲的}定理,
要找还得找\textbf{配得上它的}父亲;而只有等到那“真正的”父名——按常理本就该属于他的那
一份(而且十二年前就该……)——成为举世公认、不容置疑的共识,它才够资格到光天化日之下
露面。那篇“反常”文章,同名讨论会的一颗明珠,正是朝这个方向立下的第一块界石,由头号当
事人以他惯常的娴熟手法亲手砌下。

我的感觉是:Beilinson和Bernstein多半还挺受用——想不到竟平白沾上了对那些层的父亲名分
(那些层也被称作“反常”——同样是叫错了的),而且绑在一起的还有一位至今这般显赫的黑帮
老大——可实际上,他们是被Deligne当了枪使,好在“万一……”之际给他充当“不在场证明”。
文章既然是这样写的,任何一位消息不够灵通的读者都只能认定:那件天赐“对应”的作者,不是
别人,当然是Deligne——尽管它从头到尾没被点过名,也没被明白说出过(反正人人都该早已知道
它……)。

剩下的,只有那一点暧昧的暗影(精心算计过的),藏在一句妙笔里:那条未被点名的“关系”,
“本该在这几篇笔记里找到自己的位置……”\fnote{810}{(*)}{关于这一点,见前已引注“变戏法者”(n\textsuperscript{o}~75$'$)。}。这正是“叫停”式的笔法:含蓄然而清楚地点拨一句——却又不
肯一字一句挑明——上述那条“关系”(既然通篇并无相反的说明),至少要归到那篇光辉文章的
\textbf{三位作者之一}名下,或者(退一万步讲)由三人共享。但同样明白的是:等到时机一到
(对那位懂得等待的人而言……),无论Beilinson还是Bernstein,都不会去向Deligne争一顶实际上已
然到手的父亲名分。想必是做成过一笔交易\fnote{811}{(**)}{会想到这样一桩“交易”,是联想到了两个类似的局面。其一是Deligne与Verdier之间的交易(或许是心照不宣的,但痕迹显然可见):后者“牺牲”掉 Lefschetz--Verdier 公式,把它划进盈亏账,以成全“SGA~4$\frac{1}{2}$--SGA~5”那场手术;作为交换,他却“收进”了过世者全部的“对偶”遗产,还外带导出范畴(一篇打折的文章)作添头。(详细的来龙去脉,见“公式”那组子注,n\textsuperscript{o}~169$_5$--169$_9$。)其二是Deligne跟一位被宣告已故的大师做成的那桩“交易”——反正此人早已从流通中销声匿迹,不愁他起而反应——事涉1967/69两个年头里合开的 SGA~7 讨论班:三年之后,这个讨论班被“对半分”掉,一半归那位已故者,另一半归Deligne和一位临时拉来的搭伙人。(详情例如见注“步步升级的插曲”(n\textsuperscript{o}~169~(iii)),插曲2.)这又让人联想到跟同一位已故者(他还蒙在鼓里)就那条(据MacPherson(Mac~Pherson)说)称作“Deligne--Grothendieck”的猜想做成的“交易”(见同一条前引注的插曲1):头一半归“邮差”Deligne——是他把一条猜想(此前一直由我那班上同调学生悉心保密)通风报信给了MacPherson;后一半归那位已故者,身份则是前者的“合作者”……}——\dmark\origpage{1077}即便不曾形诸明文,至少也彼此心照:Beilinson和Bernstein分得 Kazhdan--Lusztig 猜想,外加(为
了凑足分量——反正Brylinski-Kashiwara已经
% ---------------- p.820 ----------------
趴在上面了)对那些层的共同父亲名分——那些层,我猜想是出于众人的合议,才被叫作“反常”
的\fnote{812}{(*)}{见注“反常”(n\textsuperscript{o}~76)。};至于Deligne,则分得那条没有名字的著名“关
系”,只待那指日可待的一天到来——届时他的谦逊无须劳驾挪动一步——全世界都会管它叫“Deligne定理”。而这位未来的“父亲”鼻子也着实够灵,至少把这一点嗅了个分明:就这个孩子而论
(当初他宁可拒之门外,也不肯亲自动手把他生下来……),自己做成的是一笔“好买卖”
\fnote{813}{(**)}{这笔“好买卖”,我看同时又是一桩糟透的买卖;而且即便(甚至尤其……)在事事都遂了当事人心愿的情形下,也是如此:把宝贵的天赋与创造力,挥霍在扮黑帮的营生上。}。至于
Kashiwara,他的戏份儿已经演完:在那篇光辉文章里,就那件天赐的“关系”而论,再没有人提起他,正
如再没有人提起那位当值无名之辈。人多势众围攻一个,只要那是个无名小卒,好说;可地盘上一
旦清走了闯入者,那就各人顾各人……

% ---- 以下“(f)”节, 仍属原书 p.820 ----
\subsubsection*{(f) 演员列队——或黑手党}
\tocsubsub{(f) 演员列队——或黑手党}

那份才翻开三个星期的“家庭相册”\fnote{814}{(***)}{见3月22日那条同名注(n\textsuperscript{o}~173)。},想不到又添了几
张新面孔。看得出来,“这个家”丁口大为兴旺,我这副老朽之躯都快认不过来了,更何况时世也
已大不相同。这一回,按\dmark\origpage{1078}出场的先后,是\textbf{Kashiwara}、\textbf{Hotta}
\fnote{815}{(*)}{细心的读者也许会奇怪,在这份“演员列队”(在围绕Zoghman~Mebkhout著作的那场骗局兼蒙骗之中)里,怎么见不到Kawai的名字——他是Kashiwara那篇屡被引用的文章的合作者,其第4节肆无忌惮地劫掠了Mebkhout博士论文的第三章。(关于这一点,见注“五张照片(晶体与D-模)”(n\textsuperscript{o}~171(ix)),特别是第1005页。)Mebkhout一口咬定,不能把Kawai同Kashiwara装进“同一个口袋”(Kawai嘛,闭着眼睛跟着Kashiwara走也就是了……)。他向我形容Kawai是个有点儿跟不上趟的家伙,我有个印象,他几乎已经喜欢上这个Kawai了——总之一句话,这可是他的“好日本人”,在他那里,我去碰是万万使不得的!大概也正因如此,他才没有给Kawai写信(就像他给Kashiwara的另一位搭档Hotta写过信那样),去点破他与Kashiwara合写文章里的那些骗局,从而逼他公开表态,同他的搭档兼老板团结一致。}、
\textbf{Brylinski},还有Brylinski-Kashiwara投给《数学创造》的那篇文章的
\textbf{匿名审稿人}。一帮“硬茬儿”,这准保没错,反应招数训练有素;更何况只要事关坑害一
个无名路人,只消幕后大头目一个隐而不显的暗号,便个个心领神会、行动一致(甚至等不及暗号
……)。

我又一次认出了那副\textbf{黑手党}的做派\fnote{816}{(**)}{这种离奇的感觉,去年在注“讨论会”(n\textsuperscript{o}~75$'$)(猜得出是哪一场……)里就已攫住过我:那里一派收保护费的气氛,叫人恍然如梦,或者以为正“在看一部讲黑手党称霸某座遥远大都会阴暗底层的电影……”。这一回,在这趟穿行于当值无名之辈种种倒霉遭遇的漫游里,这种感觉又寸步不离地陪伴着我……}:在他们无可争
议的领地上称王称霸,而领地的核心,便是代数簇以及诸如此类的上同调理论。一帮才华耀眼而心
狠手辣的人物,脑瓜子个个无懈可击;在所谓“当值无名之辈”那场行动接连四段的插曲里,我一
路重新见识了他们的身手,直到“反常层会议”登上顶峰。除我方才点到的四位黑帮老大(其中一
位匿名)之外,还请诸君记起“硬核”的另外五名成员;这样一共九位,为埋葬\textbf{那位闯入
者,那个不是他们一伙的人}而全体出动。

有大头目\textbf{Pierre~Deligne}——那位永远懂得“下水”下得最少、进账却捞得最多的主儿。
有他的副手\textbf{Jean-Louis~Verdier},人称“大善人”——正是此公,主持过某位无名之辈某
篇论文的答辩委员会;又正是此公,当过那场令人难忘的讨论会的两位组织者之一,那场会肆无忌
惮地掠夺了同一位无名之辈。另一位主其事的是\textbf{Teissier},他同Verdier联名为那场令
人难忘的讨论会的令人难忘的文集,签下了那篇令人难忘的引言。跟旁人不同,他\dmark\origpage{1079}似乎不过
是跑了个龙套、当了个挂名替身,自己从中一无所得——唯一的乐趣,不过是讨得
他知道既有名望又毫无顾忌的那班人的欢心。最后\fnote{817}{(*)}{(5月25日)这个“最后”终究是言之过早——打那以后,又有别的帮派成员自行冒了头,进入了我的视线。关于这一点,见“光荣之日”注(n\textsuperscript{o}~171(iv))里打字稿第962页的脚注(*)。\par(5月30日)最新消息:又一名成员——R.~Remmert——刚刚被认了出来。关于这一点,见本注 c) 部分(“记忆的失灵——或新编历史”)。}还有\textbf{Beilinson}和\textbf{Bernstein}(对这两
位,就在这几页里
% ---------------- p.821 ----------------
我刚刚有过一番更充分的领教),被一根根看不见的提线轻轻地牵动着……

我还会等着,不急不躁,也不抱幻想,看未来还会给我们预备下哪些别的“反常层会议”,有整
个教团毫无保留的首肯作保,好为“科学”增最大的光荣,为“人类精神的荣誉”添彩。

\subsubsection*{c3.~根与孤独}
\tocsubsub{c3.~根与孤独}

\notehead{171$_{3}$}{}

(4月18日)一步一步跟完我朋友Zoghman那串倒霉遭遇的这第四天行将结束之际,我比去年更懂得了一些态度与心性——尤
其是冲着我来的那些——而直到去年,它们在我眼里还显得古怪。总而言之,凭着他那份自知分量的工作,他原以为进了“一个大家庭”,
有点像那位从不被人提起的已故宗师的家——的确,无人提起;可即便不提,他分明在场。谁承想,他
发现自己落进了一个鲨鱼的世界:个个彬彬有礼、甚至和蔼可亲,牙齿却毫不留情——转眼之间,他
带来的东西被剥夺得一干二净:八年孤独劳作的果实;然后有人让他明白:你,我们已经看够了——一个讨厌
的、不请自来的外人。换了他那个位置,没几个人能不留下创伤。我不知道他可曾对什么人敞开过他
的挫败,除了几声苦涩的暗示,又含混得反倒像在作证对他不利:一个怨气冲天的人,边缘上还有点不合群。

尽管没有人点我的名,我却俨然是这伙毫无顾忌、毫不留情的世界的“爸爸”;他实在没有什么理由
信任我。诚然,我们的初次见面——1980年,那时他还远在千里之外,猜不到等着他的是什么——曾
经打下过信任的底子;我分明感到,这个底子,顶着一切,一直保存到今天。说到底,他心里明白:我
纵然是一切鲨鱼的“爸爸”,也不会像他们那样行事。可\textbf{怨怼}是有的,这确定无疑;而且它乐于\dmark{}\origpage{1080}
摆出一副打心底不信任的架势——虽然(至少我感觉如此)那是“贴上去的”。

为自己认定的正当权益“斗争”,这容易——只要你属于一个团体,多小都行,同它声气相通。可那个
孤身对抗全世界的人、那个被排斥者、那个不请自来的外人,就像一棵被夺走了沃土的树。他身上的
力量帮不了他;那力量变成苦涩,反过来咬他自己,仿佛要同那个排斥他的全世界合唱一个调子。

当我手里捧起那本书——它为 Motif 掘墓复生,同时埋葬了让它们现形的工匠;那本书,署着一代
光辉人物(其中有我参与培养的)里最光彩的四位作者的大名——当我终于读到它(纯属偶然,因为
此前没有任何人觉得有什么值得特别提醒我的……),那一刻,在我结识数学家世界三十六年之后,
我第一次知道:\textbf{我孤身一人,对抗所有人}。过去八年里发生的许多事,忽然拼合起来,显出全部
意义。当一个人蓦地重新发现这份孤独,那滋味很是古怪。那一天我不得不缓一口气;此后一连几个
星期,一天一天地认清\emph{葬礼}的全部规模——一场与那部著作相称的\emph{葬礼}。

可这同Zoghman毫无共同之处——他的“自己人”在他还没真正扎下根之前,就把他“挂账”了。而命
运对我,是笑过的。多亏长者们接待了我(他们后来故去也罢、退休也罢、早就不再搞数学也罢,统
统不要紧)——多亏我年轻岁月里寻得的兄弟般的接待,我,得以在我自己选定的土壤里,“扎下根
来”。这些根,扎下去,长起来,随着岁月,变得深沉而有力。这些根,牢牢栽在一片不是“共识”、也
不是任何时尚的土壤里——扎得想必比任何以制造时尚、追随时尚为乐的人都深\fnote{818}{(*)}{如果说我在数学里也好、别处也好,从来不曾操心去追随或制造时尚,我知道,这正是我有幸在幼年便得以发展起来的那种深根的表现之一。既然我从一开始就在自身之中有了深根,投进我的大事业里的那份精力,就不会被补偿性的饥渴——比如领风气之先的饥渴,或保持并显得合乎必备之“调门”的饥渴——所分散。
关于我的童年和这些“根”(我想,我没有说出这个词),我在注“天真(阴阳的联姻)”(n\textsuperscript{o}~107)里说得具体。}。

\pmark\origpage{1081}说到底,我还“孤身对抗所有人”得起——想说什么就说什么,走我自己的路。

(5月25日)\fnote{819}{(*)}{下面两页,原本是作为注“……与意外之喜”(n\textsuperscript{o}~171(iii))的一条脚注写的。在哪里安插它们,我犹豫了一阵,最终决定收进眼下这条“根与孤独”里。因为这是“封神”全部注中唯一的例外:我试着从自己的经历出发,尽力去体会Zoghman本人如何经历了那些我为之作编年史的事件与处境。}不消多少想象力,就能体会Mebkhout的沮丧:他的核心结果之力一旦得到承认,他
就感到自己像一根麦秸似的被“扫掉”了\fnote{820}{(*)}{“扫掉”一语,借自Mebkhout的一封信(比正文所引那封早两天),相关段落我在这里照录:
“的确,[Kashiwara的]那个可构造性定理……让我爆发了。而且从那一刻起,像Deligne这样的人,一眨眼就
能找到我全部的结果,连‘好上帝定理’的各种形态在内,证明嘛,照你的说法,四个小菜的功夫。这就
说明了为什么这一切几天之内就被扫光了。”
我想,Mebkhout在这里把一个Deligne式的、默而不宣的“推理”,说得再确切不过:据他人\textbf{本来可
以}(而且\textbf{本应})凭他自己的手段、行囊与一切,“四个小菜的功夫”找到这些果实,所以就可以据为己有。这类推理——在类似的情境里,人们常常禁不住要做——唯一的别扭之处在于:\textbf{难的全是想
到}。而真正“想到”了的,是Mebkhout,绝不是Deligne,更不是别的什么人。创造不归“技术”管:技术,
在终于有人看见了无人曾看见的东西之后,顷刻间便可“扫光”一个局面,快过把它写下来。创造不在
“扫光”里,而在\textbf{看见}无人曾看见的东西之举里;用自己的眼睛看,不“追随”任何人。把这两者区分
开来——把创造之举,同摇一圈就转圆的那根曲柄区分开来——正是数学家操守的一部分。}。他给我写信(4月24日,在他新近来我家之后)说:“我花了八年,搭起
Kazhdan--Lusztig 证明所用的那些结果。他们用一个星期就把它证明了。”这一回,还是有一种羞
怯拉住了他,不肯把心里真真切切感到的那些话说到底;\dmark{}\origpage{1082}我在这里替他补上那句“没说出的”:而事情一旦做
成,“他们”便拿着那件别人在孤独里打造出来的崭新工具,彼此之间昂然招摇,同时让工匠明白:你,
我们已经看够了……

可这件事实在大得离谱,以至于当场,Zoghman还不大相信他自己那健全官能的见证——正如去年5
月2日,读到吕米尼讨论会文集时,我自己也迟迟难以相信自己官能的见证\fnote{821}{(*)}{关于这场讨论会(1981年6月),见注“不义——或归来的意义”或“光荣之日”(n\textsuperscript{o}s~75、171(iv))。说实话,去年5月头一个星期写下“送葬行列之七:讨论会——或Mebkhout层与反常”(n\textsuperscript{o}s~75--80),还不足以克服那近乎不可克服的惰性——在一个严格说来只有你一人使用这种见识的处境里,“相信自己健全官能的见证”何其难。直到五个月之后,终于“有血有肉地”面对了现实——我的朋友皮埃尔(Deligne)来到我隐居处看我——那股隐秘而顽强的怀疑,才终于消散。关于这一点,见注“义务已尽——或真话时刻”(n\textsuperscript{o}~163),尤其打字稿第782--784页。}。Zoghman是在去年1月读到这同一部
文集时——在 Kazhdan--Lusztig“总彩排”三年之后——才终于多少弄明白,实际发生的是什么事。

打击是可怕的,我自以为懂得——Zoghman当场以为他要把命搭上。幸好人这身子很结实——Zoghman
今天还活着,其间还结了婚,做了父亲……可我想,即便手捧“文集”的那一刻,他仍旧不能完全相
信。总有什么东西“卡着”。说不定,在我写这些行的时候,他仍旧没有全信。应当说,即便只按“理性”
或“客观”的字面讲,这件事也实在难以置信、实在巨大:直到今天,除了我(或许除了他,也未必
……),竟没有\textbf{一个人}敢相信自己眼睛、看见它——虽然它比一座主教堂还大!

可对一个被当头击中的人——被他的那些敬重的、什么都不缺的长辈,施以犬儒而\textbf{无缘无故}的
不义——这种事,只怕永远没法全信,只怕永远“\textbf{超出理解}”……而恰恰是这种事,能够毁掉一个
人的一生。它们之所以有这份\dmark{}\origpage{1083}毁灭之力,是因为一种幽暗的、绝望地被压下、却无可否认的知觉:
\textbf{毁掉}的\textbf{意图}——就这样无缘无故地,为了\textbf{取乐}:为了取乐于漫不经心的一挥手,碾碎那
对你珍贵的东西,碾碎(如果办得到)那构成你生命的实质与盐分的东西。正是这种“无缘无故”的
恶意里的反常快乐,才真正“超出理解”……

我想,Zoghman从来不曾真正对什么人谈过这些,大事之前也罢,之后也罢——至多几个单音节的
词,除了他自己,谁也解不开。单是 Kazhdan--Lusztig 那一幕,已经太大、太不像话,他不敢指望谁
会信。根深蒂固的共识,把最明显、最明白、最无可抵赖的事实,像麦秸一样扫掉。而这里说的,
是一件疼得那么切近、在他的存在里那么“一碰就疼”的事:只消想到,他若向谁袒露,那人可能拒斥这
份不速的信息,他面对“超出理解之物”的痛苦可能不被接纳——这个风险、这个概率,便膨胀成
\textbf{不堪忍受}之物;这种事,付任何代价也不去碰——哪怕为此当场憋死,要憋就憋……

两年前,他倒是对我“用单音节的词”说过。说不定,在他心底,他盼着我会听懂这些单音节——不只
按字面,而是同时听出他不敢高声说出的一切(也许连对他自己都不敢说……)。这盼望确实疯得
可以(在一个一切都疯得没边的处境里!);对我后来确凿得知的那些,我当年连想象的边都没沾
到。只能如此,谁叫我那时没有一份详尽的、原原本本的信息呢\fnote{822}{(*)}{(6月1日)更公平的说法或许是:以我这种有限的敞开与在场的状态,除非在极罕见的时刻,事情“不可能是别样”。然而我相信,我们人人都有“耳中之耳”,完全听得见没有说出的那部分——只不过,我们通常留神把这只耳朵捕到的讯息,排除在自觉注意的领域之外……}。而Zoghman那一边,也远不敢把这份信息
给我。这很疯;可这并没有妨碍他怨我。他总得怨一个人,怨一个足够亲近的、摸得着的人,好把“超
出理解之物”在他心里搅起的东西,至少卸下一部分在他身上,好让自己从那啃噬着他的东西里,解
脱那么一点点。

\subsubsection*{c4.~掠夺的空白授权——或“高尚行当”}
\tocsubsub{c4.~掠夺的空白授权——或“高尚行当”}

\notehead{171$_{4}$}{}
\dmark{}\origpage{1084}(6月2日)我给\emph{葬礼}画上“句号”、连同末尾那条注“自深处”(4月7日)一道的那份心满意足,到今天已满两个月——而这两个月来,我也一直在拼着老命,给\emph{葬礼}的最后一部分“做最后的润色”!这几乎是去年同一时节那出戏的原样重演——当时,我给那即将成为\emph{葬礼}第一部分的东西“收尾”,也是收个没完。
% ---------------- p.823 ----------------
和眼下一样,拖个没完的正是“最后一刻”——拖到我连吃喝都顾不上,最顾不上的还是睡觉。就这么一直拖下去,直到我的身体油尽灯枯,宣布弃权。那正是整整一年之前(差不了几天),我不得不撂下一切,足足三个多月,一门心思只顾从一场急性衰竭里爬出来\fnote{823}{(*)}{关于这一段,见注“事件——或身体与精神”(n\textsuperscript{o}~98)。}。可这一回我多了个心眼,格外当心,不再走同一条老路。这条命,我还想留着呢……

这一回,又是那场没完没了、一弹再弹的“调查”。我原先估计,撑死也就一条十来页的注,取名“四种运算”,把去年那阵旋风式调查的结果“归置归置”、做个总结。可如今,调查卷土重来、干得热火朝天,眼看要满四个月了;十页变成了三百页上下,而且这还不算(完全)完事!我再不敢作什么预言——从9月底重新动工算起,已经到第九个月了,我一直在“马上就要完工”!只有等到最后一批笔记全部打成清稿、通读修改完毕、交去复制的那一天,我才会知道它\textbf{真的}完成了。(到那之后,剩下的就不再是我的活了。)我唯一知道的,是我巴不得赶紧到那一步,就像巴不得看到一场漫长而磨人的疾病有个尽头;而我必须尽我所能走到终点,不被种种想象出来的期限推搡着走。只有等该\textbf{在此时}看见和说出的都已看见、都已说出,我才会停下来喘一口气。

害我吃了最多苦头的,正是那该死的“\emph{封神}”——我说不上为什么。这几样“四种运算”,是《收获与播种》全书里唯一一处写得磕磕绊绊、零敲碎打、一路费力挣扎的部分——按说本该是现成的活儿:不过是一次简单的“归置”,没错;没有任何东西会以触到\dmark{}\origpage{1085}要害的方式,把我这个人牵连进去、拿我这个人问罪,从而动员起抵抗的力量,一个“\textbf{摩擦}”。可天知道这摩擦有多少,而且在\emph{封神}这里,比其余一切加起来还多!它打哪儿来的?

早在“操演”那条注里,写起来就已经够费劲了。没完没了地拉长,就是从那里开始的。光这一道“运算”,最后就足足写了八十页——而如今,一个月之后,\emph{封神}已经写出了整整两倍的篇幅。话虽如此,除了“操演”里兴许有那么几页(边上带点儿太浓的“侦探”味儿……)(在那里,我兴许一头扎进一桩离谱透顶的“骗局”的絮叨细节里,扎得比必要的还深……)——除了这件写得眉目详尽、对“不在局中”的读者来说多半还有点磨人的“凭卷宗办案”的活儿之外,我并不觉得我最终码出来的那一摞摞上百页的东西是画蛇添足,是回锅的冷饭,是头发丝里切四瓣的功夫。让我一直提着气撑下来的,恰恰是\textbf{新的实质内容}大量地、出人意料地涌到我身上,它们我无论如何都得安置妥当,由不得我愿不愿意——其中甚至有,可不是嘛,数学的实质内容!有好些时刻我几乎招架不住:一时要落到白纸黑字上的东西实在太多,而且刻不容缓——这些事情一个个都还热乎着,甚至烫手,可你偏偏只能一件挨一件地去料理……

这样的丰富,本身就自带着一股强劲的工作推动力,丝毫不是那种会惹出“摩擦”的东西,恰恰相反。这摩擦,可以肯定,并非来自实质内容本身,而是来自我投进这份工作里的那股唯我的执念之力。说来像是悖论:恰恰是我那股“赶紧收场”、“一把甩到桌面上”的不耐——急着说出我该说的这样那样的事,而它们就发生在此时此刻,同我息息相关、触我至深——正是这种不耐(我想),制造了摩擦,制造了能量的涣散。摩擦是分裂的征兆,是几股力量朝相反的方向对拉,每一股都因对方的抵抗而愈激愈怒:
% ---------------- p.824 ----------------
一头是“赶紧收场”的急切——趁我还在打磨这块料,就想把它“甩手”交出去;另一头,是要求把眼下这一刻让我隐约望见的东西走到底,不满足于差不离,不让人推搡,也不把自己关进一个赶着“结项”的“计划”、一份预先定死的“日程”里。我很清楚:一旦把“意外”排除在外——这个防着我原地打转的搅局鬼——我的工作就失去了品质和意义。它\dmark{}\origpage{1086}就沦为“文牍差事”。这些年下来,我对这处看似不算什么、实则是一切的“小小差别”,已经变得非常敏感。这样的大转弯,如今偶尔还会启动,总在沉重已极的时刻——但从来没有维持多久。一旦走上这条路,那孩子把一切都扫到一边去——再想撑下去也是白搭。工作本身的兴致,那份\textbf{渴望}——它全然不同于攒页数的饿劲儿、不同于抢着画句号的急色——霎时间无影无踪,你便傻乎乎地守着涂黑纸张的活计。到那时,真就再没意思了——剩下的只有一件事:把航向拨正,而且马上就拨!

工作里总有一种\textbf{不耐}(我的一个老相识……),不停地把我往前拽。我觉得,这和自我跟这几样“四种运算”较上劲以来沉沉压在我身上的那种不耐,不是同一种。那另一种不耐,不是压下来的重量,而分明是一种往前拽的力量。它是胃口好的征兆,不是倦怠、疲乏或餍足的征兆。它不是累积的不耐,不是收尾的不耐,不是赶着把一份“计划”“结项”的不耐,而是要认识我面前那个即将向我把底交出的未知者的不耐。是一个赤身的孩子,独自站在无边的大海面前,急于纵身潜入其中去认识它……\fnote{824}{(*)}{这一形象,已见于注“孩子与海——或信与疑”(n\textsuperscript{o}~103)。}

不过,现在该回到我朋友Zoghman那串倒霉遭遇的叙述上来了——就在这条预定作为\emph{封神}最末一笔的注里。前面说过,这段叙述,Zoghman本人从来只是东一块西一块地零星递到我手里:随着信件、电话、会面,碰上什么算什么;可以肯定,我这边思考的推进和\emph{葬礼}的写作都因此受了影响,至少在专写我朋友那番坎坷的那一部分里是如此。如今我更懂得这种欲言又止的意味了,因为对“受害者”角色的一切执着(去年我还自以为从中看出几分)早已烟消云散(就算它当初确实在过)。我这边,想必也有过那么一些时刻,多少有些饱和了,态度上是那么一副“求您别再倒了,行行好!”的架势。这八成也没给他鼓劲儿。说实在的,那句“这儿是日本人、那儿是Kashiwara”的老调,Zoghman想必唱了总有四五年了,叫我不胜其烦;他跟他们之间确实领教过一些事,这不假。可我心里明白:他之所以领教过那些事,他的成果之所以就这样被交出去任人掠夺——近乎官方式地:“来吧,好人们,放量地拿,千万别客气……!”——那\textbf{不是}因为几个遥远\dmark{}\origpage{1087}的日本人。这\textbf{是因为“自己人”}:是“小家庭”\fnote{825}{(*)}{(6月16日)关于这一点,Mebkhout特意要我强调:他早已完全不再认同于这里所说的那个“小家庭”。}里的人——都是我们自家阵营里的体面人物,而他从不点名,除非引用他们的著作,并带着配得上他们崇高声望的全部敬意。

至于Kashiwara之流,我可再也不想听人提起了!眼看事情卡住了,Zoghman这时拿出他的智慧和耐心,索性撂开不谈,却并没有因此放下对我工作的关切,也没有停止在这里那里默默地、实实在在地帮我一把。

直到4月初他最后一次来我这儿,我才终于看了看那个“日本包裹”。起初还有点勉强。我原以为,自己准会陷在一堆解不开的超技术纠葛和读不懂的文章里腻烦得要死(没准儿还是日文的……),反正我是永远不会去读的——结果不然!事情简单得像打声招呼——有点像巴黎地铁(或者该说东京地铁)里“扒手的故事”。说真的,还挺有趣(至少,被摸走钱包的是别人的时候……)。
% ---------------- p.825 ----------------
这一来,Zoghman和我之间的僵局也松动了,他的倒霉史,这里那里、一道道闪电似的,也让我听到了些零头碎块。原先我多少按“技术情报卡片”的样式记下的几段插曲,如今被鲜活的临场追忆充实起来;恰恰是这一类东西,在科学文本那副纹丝不动的“立正”姿态里,甚至在同事之间的通信里,似乎永远被禁绝——您总不会想看吧!在“四种运算”里,我甚至不得不使劲摇晃自己,免得恰恰又落回那种腔调——“侦查结论”式的腔调(甚至,“诉苦清单”式的……)。Zoghman交出来的这些“碎块”,帮我从里面走了出来,帮我同一种活生生的实质保持接触。

Zoghman离我家的当天,我就重新拾起\emph{封神}的活儿,趁他讲给我听的东西还没凉,再写它一条两条子注。这就有了那几条注(或者说子注——写得多了,我自己也数不清了……):“眼光的绽放——或闯入者”、“黑手党”(后来我把它又分成七部分,每部分起了名字),还有“根与孤独”。我把整个儿火速寄给了他,请他在我送去打字之前提提意见。到这一步,我觉得自己多少有点在替他说话,我要确证:凡是我根据他讲的转述出来的东西,都得到他毫无保留的认可。他随即寄回了详细的意见(4月22日与24日的信)。

\pmark\origpage{1088}这些意见里,不乏这样的“碎块”,为我那些注里偶尔显得有点瘦骨嶙峋的事实骨架,敷上了鲜活的血肉。

也正是这样,我才知道Zoghman当时在场——就在那个令人难忘的1980年4月22日,Goulaouic--Schwartz讨论班上。就在那一天,Kashiwara把“好上帝定理”当成自家出产的定理当众宣布——而这条定理,是他几个月前在莱苏什讨论会上从Mebkhout亲口听来的\fnote{826}{(*)}{关于莱苏什讨论会与Goulaouic--Schwartz讨论班上的这段插曲,见注“黑手党”(n\textsuperscript{o}~171),部分(b)“最初的麻烦——或太平洋彼岸的黑帮老大”。}!这胆子实在大得没边,何况Mebkhout人就坐在底下,难怪听起来像天方夜谭。Mebkhout当场没有爆发(我直纳闷他是怎么忍住的……)。他客客气气地等报告讲完,“为的是公开抗议这种做派,向他提起莱苏什的那场会议、提起他的那个问题\fnote{827}{(**)}{指的是1979年9月莱苏什讨论会上,Mebkhout报告结束时Kashiwara所提的那个问题。关于此事,见上一条脚注中所引的那条注。}。Goulaouic请我私下解决我们之间的事。大厅几秒钟之内就空了”。

这便是那些“碎块”中的一块,由这段惜墨如金的叙述交了出来。后来我又在电话里听到些细节。这件事值得停下来细看。关于八十年代数学界的风气民情,它透露的东西可不少。这里涉及的,已不是某位心狠牙长的“黑帮老大”的心态——科学界传统价值败坏的一种极端症候——甚至也不是那班体面风光、处处周全的“建制中人”——在他们身上起作用的,是护着“自己人”的阶级本能。在这里,是整个大厅眨眼之间一空——转眼间一个人都不剩\fnote{828}{(***)}{这幅景象在我心头不可抗拒地勾起一连串联想,同三年前我亲历的一个场面如出一辙:那是一场Bourbaki讨论班的末尾,承人们好意赏我十分钟,让我谈谈某条坑害外国人的凶恶法令。关于此事,见“我的告别,或:外国人”一节,n\textsuperscript{o}~24。}!你们自己私了吧——我们嘛,什么也不想打听……

我真想知道,Goulaouic和那场讨论班上其他安安稳稳的听众,脑子里当时到底转过了什么念头——在那儿做报告的,可是一位声望斐然的外国讲演人(讲的题目,我相信他们当中没有谁真正熟悉)。这场插曲,说到底,本来够人好好想一想。但我\dmark{}\origpage{1089}恐怕他们谁也没费这份心,我更愿意猜想:全体一致,大家争先恐后把这段难堪的插曲忘了个干净。不过说到底,只要肯费点心思想一想,而不是拔腿就逃,这桩阴暗的事情里,毕竟还是有\textbf{一件}事是清楚的。Mebkhout的语气和措辞(何况他们认识此人,至少在各处讨论班里同他打过照面),使人几乎不容怀疑:
% ---------------- p.826 ----------------
这件事里必有\textbf{一个骗子}——不是Mebkhout,就是Kashiwara。

当然可能,他们在心底里早已提前断好了案:Mebkhout在编故事——那位尊贵的来访客人会去掠夺一个无名听众,这叫人怎么想象!这便等于说:面对一个无名之辈,名人无论做什么,都在一切嫌疑之上。这便是\textbf{掠夺的空白授权},交给有名望的人,用来对付求告无门的人。他要说的话,不会有人听:“你们自己私了吧!”

要不,他们就是一头埋进了怀疑里:谁知道谁说真话、谁说假话?(更何况,只要捂起耳朵,就更不知道了!)不错,一个Kashiwara,当着事主的面公然掠夺一个不知名的陌生人,这份蛮横的厚颜,看上去几乎难以置信。可说到底,还有更难置信的事:一个无名的陌生人(他们全都认识他,而他也从未靠骗子的把戏、靠那份胆量在他们眼前露过一手……)竟敢当众指控一位Kashiwara犯下粗暴的剽窃——假如他要说的纯属凭空编造……再退一步,就算他讲的或许有根有据,拿一句“你们自己私了吧!”把他打发到九霄云外,这依然是掠夺的空白授权。这就好比,一个人在大街上被一群穿燕尾服的恶棍洗劫,喊一声“抓贼呀!”——人们却冲他喊:“你们俩自己私了吧!”

况且据说,长期以来,纽约和其他美国大城市的下等街区,事情就是这样办的——在那里,谁也不想跟当地称王称霸的黑手党找麻烦。反正如今,数学界里事情就是这样办的(从什么时候起,我说不上来),在那些被当作“上等街区”的地方就是这样办的——比如Goulaouic--Schwartz讨论班\fnote{829}{(*)}{在这里我可以高兴地指明:那桩令人难忘的插曲发生的那一天,“他的”讨论班上,Laurent~Schwartz本人并不在场。后来他是否听人说起过,我就不得而知了。},或者在那班“搞”代数簇上同调的名流当中。

\pmark\origpage{1090}若按理性的说法、照字面来理解,“你们自己私了吧”这句话简直近于低能——何况无论怎样,两造之中必有一方是不老实的,这一点再清楚不过。在心理层面上,这句低能的话所表达的,是在一种被感到“碍事”的处境面前的一种\textbf{卸责}。这也是对下述明摆着事实的蓄意无视:遵守数学家这一行当基本伦理规则的问题,绝不是一件纯属“私下”的事,可以由僭取了蔑视这些规则之权的一方,同吃了亏的另一方,关起门来了结。这是一件\textbf{公共事务},一件关系到\textbf{每一个}数学家的事。

正是靠着普遍的冷漠,靠着每个人在承担个人责任时的惊慌失措,黑帮心态和“反常层会议”(Colloque~Pervers)那样无耻透顶的勾当,才能在科学界肆无忌惮地兴旺繁盛。一些人的惊慌,同另一些人的无耻,是\textbf{同一种腐败}的正反两面。1980年4月22日那一天拔腿就逃、捂着耳朵的那些人,同样促成了那场令人难忘的讨论会在次年的\emph{封神},其功劳丝毫不亚于那班凭空搭起这场辉煌大骗局、又神气活现跑去里面招摇过市的黑帮老大。

(6月3日)也正是在Mebkhout最后一次来我这里的时候,我从他那里听到了一些很有教益的细节,涉及同一场辉煌讨论会的某几位与会者,以及他们各自身上竞相盛开、一朵赛过一朵的“新风格”。我有机会翻了翻那次会议的报告集第二卷,里面有Verdier的文章、Brylinski--Malgrange的文章,又扫了一眼Laumon的博士论文(用一双比初次收到它那天更有经验、也更不分心的眼睛)。这篇论文其实是同N.~Katz合作的工作。关于这些工作里遵循的“新风格”,我在注“光荣之日”(上帝知道,这个名字它受之无愧……)的那条长长的脚注(打字稿第962页)里写了几句评论;在那条注里,为了更细的情节,我还指向了眼下这条注
% ---------------- p.827 ----------------
(当时它还没写出来)。说过的话,就要兑现!

Mebkhout给我讲过,他有怎样的荣幸和便宜,得以两次向N.~Katz谈起自己关于对偶的想法、关于连续系数与离散系数之间联系的想法。第一次是在雷恩的$p$-进分析讨论会上,1978年7月。当时他\dmark{}\origpage{1091}在“小圈子”里讲解了他的定理:复解析空间上D-模的整体对偶定理——那条为Serre对偶和Poincaré对偶加冕的定理\fnote{830}{(*)}{关于这条定理,注“作品……”与“三座路标——或天真”(n\textsuperscript{o}~171(ii)、(x))里都有谈到。}。在场的有Katz,有Illusie——正是在\emph{葬礼}里已经不止一次谈到过的那两位。Illusie呢,像他惯常那样和和气气、客客气气,觉得这真是漂亮极了——大致是这么句话\fnote{831}{(**)}{况且这本身还是一种分文不费的“客气”。两人的反应路数彼此不同(Illusie是“阴”,Katz是“阳”),实质却并无二致:只要是“无名先生”(Monsieur~Personne)说出来的,那就从这只耳朵进、打那只耳朵出!关于这一点,见注“蒙骗”(n\textsuperscript{o}~85$'$),尤其是打字稿第351页上我对Illusie的那几句观察。}。至于Katz——我猜,他在那儿是平生头一回听人谈起D-模(那阵子,这东西还远不是那场令人难忘的讨论会之后的那种大时髦)——他只干巴巴地撂下一句“这早就有人知道了!”,随即转身就走。既然跟他说话的,不过是个不知打哪儿冒出来的“无名先生”——而他N.~Katz,就在同一年,还要在几千位尊贵同行面前发表演说,向新科菲尔兹奖得主Pierre~Deligne致敬呢……——那当然只能“早就有人知道了”。

第二次,是1979年9月莱苏什讨论会之后不久\fnote{832}{(***)}{关于莱苏什讨论会,以及Kashiwara在Goulaouic--Schwartz讨论班上的那桩骗术,见注“黑手党”(n\textsuperscript{o}~171),部分(b),“最初的麻烦——或太平洋彼岸的黑帮老大”。}。Katz那时在法国高等科学研究所(IHÉS)。Katz在$p$-进微分系统方面的功力尽人皆知,而Mebkhout心里明白,这同他刚在莱苏什讲过的“好上帝定理”多少有些干系;于是他专程赶到法国高等科学研究所,把他在莱苏什的文章送上门去,同Katz谈自己的想法和结果。有了雷恩那场接待打底,可以说他这人是有点韧劲的,竟不知疲倦!反正,这一回的戏码跟上一回差不多。Katz又一次居高临下地接待了这个不知打哪儿来的无名之辈——此人竟敢第二次跑来催问,而且没准儿连招呼都没事先打。人一旦成了大人物,有时就不知道该怎么躲开这些缠磨人的家伙了……

一年之后,只消这同一批想法——由一个无名的陌生人,在孤独中久久揣着、慢慢养熟——跟着那场Katz先生遗憾未能赏光的辉煌讨论会的余威,被四处吹得震天响,说成是某位Deligne(或者某位Kashiwara,已经说不清了……)最新挖到的宝贝,它们便\dmark{}\origpage{1092}在这位大人物眼里陡然兼得了分量与斤两。替他把来龙去脉讲解清楚的,准是Laumon——Deligne最出众的弟子之一。而这位Laumon,对这些想法的来历也一清二楚,而且是第一手:因为正是那个无名的陌生人亲口告诉过他。可弟子以踏着老师的足迹为荣;而老师呢,早已明明白白、毫不含糊地示范过:对那个被判给沉默与湮没的人,该采取什么样的做派。

台前灯火的荣耀,归Deligne们,归Verdier们;Brylinski、Katz、Laumon诸公,也都赶在好时辰跑来,好从中分得自己的一份!乐声与吹吹打打归他们,感恩戴德的群众的欢呼也归他们——这群人欢天喜地蜂拥而至,庆祝这落入“新主子”们之手的“高尚行当”。

\subsubsection*{墓后尾声——或洗劫}
\tocsubsub{墓后尾声——或洗劫}

\notehead{171$'$}{}

\pmark\origpage{1093}(6月14日)直到一个月以前,我始终觉得,“葬礼”的精神所及,仅限于我间或称作数学“上流社会”或“大世界”的那个天地,尤其是那个天地里我素来出入、自己也身列其中的那些圈子。在朗格多克科学技术大学(USTL,蒙彼利埃)——它十二年来一直是我的挂靠机构——我不曾察觉到任何排斥的迹象,任何故示轻慢、失礼乃至粗暴的姿态,与那场十五年来方兴未艾的“葬礼”同一走向\fnote{833}{(*)}{在注93(第396页第3段)里,我尤其表达的就是这层意思。}。一桩新的事实刚刚闯入这幅安宁的图画,以一种严酷的方式改写了这幅图画,也改写了我同我的挂靠机构之间的关系。

% ---------------- p.828 ----------------
依照积习已深的机制,我起先并没有想到要把这桩新近的插曲收进我的见证《收获与播种》——乍一看,它出现在这里,仿佛“汤里落进了头发”,好不是地方。是顶住了几分切实的抗拒之后,我才终于承认:把这段插曲按下不表,便是有负于我的见证的精神。诚然,这段插曲余温犹在,何况我这一次“承受”得相当吃痛——这倒又给我方才影射的那些“积习已深的机制”平添了一分力气。可这一次,我领教这桩事故那些雄辩而不合时宜的教训时,感受如此鲜活,这本身也是一个征兆:它触及我至深——触及我的职业活动,触及我同我所隶属的职业圈子的纽带。所以,这恰恰属于《收获与播种》立志要作深入见证的那类事情:不允许有任何“保留的角落”,是我禁绝自己去触碰的——无论那“谨慎”是对我自己,还是对任何人,都用错了地方。

此外,在我对“葬礼”的思考这个更专门的框架之内,我感到一目了然:两者之间有着直接的关联。有可能,这类关联并不是单纯的因果关系:当地的某些同事,或许是终于把“葬礼”当真看在眼里,于是认定,从今往后他们也不妨“照样来上一场”。即便果真有这样一层因果,依我看,它触及的也只是事态一个附带的、偶然的侧面。然而另有一个更其本质的侧面,也是最令我震动的,无论在大写的“科学大世界”里,还是在一所外省的平凡大学里,皆无二致:那就是学术与大学环境中某种\pmark\origpage{1094}\textbf{退化}——也许是空前的退化:人际关系的品质在退化,礼数与尊重他人这些最起码的形式在退化,科学伦理亦然;而科学伦理,本就与尊重他人、尊重自己牢不可分。因此,下面这些页,读者不妨看作对“一个时代的风俗画”——在数学环境里,想必是一个时代的末日——的一份贡献(在围绕“葬礼”的思考一路提供的众多贡献之外)。

与其在这里把事件的经过或详或略地重述一遍,我更愿意转录四份\textbf{文件}——它们同样能把事件描述清楚。即:
\begin{enumerate}
\item 一封“致USTL数学任教同事书”,日期为5月28日,我在信中向他们通报了某种情况,并表示希望在全体大会上加以讨论;
\item 一份“答复”,出自Charles夫人——USTL数学楼的房舍负责人——形式为5月30日的一封通函,抬头具我的名,而实际上寄给全体数学教师;
\item 一项决议,由UER~5(第5教学与研究单位)全体大会于6月6日表决通过,大会议程为:“关于Grothendieck教授办公室搬迁的情况通报与讨论”;最后则是
\item 一封“致数学楼旧日同事书”,日期为次日6月7日。
\end{enumerate}

% ---------------- p.829 ----------------
我没有把两封信列入这些文件:一封是我5月21日致Charles夫人的信(第一份文件里提到它),另一封是我致USTL临时行政负责人Cano先生的信(这份文件里提到它,第四份文件——即“一场误会的尾声”——里也提到它);比起下面转录的文件所含的内容,这两封信在我看来没有提供任何新的信息。

对Charles夫人的信,我只作一条评注(“事实上很难联系上他”——这里的“他”,指的是区区在下,而这封信名义上正是写给我的):蒙彼利埃寄到我住处的信,一天即可到达;而多年以来,我离开住处,唯有去USTL的时候。

% ---------------- p.830 ----------------
\begin{quote}
\pmark\origpage{1095}朗格多克科学技术大学\\
数学研究所\\
数学楼里的一场洗劫\\
致USTL数学任教同事书\\
Alexandre~Grothendieck 撰

\hfill 蒙彼利埃,1985年5月28日

亲爱的同事:

上星期,我托UER的一位女秘书去我四楼的办公室取一份文稿,才得知那间办公室已被搬空,我的东西荡然无存——这一点我就在今天亲自核实过:只剩下一片光地板。从没有人通知我:我的办公室竟可以不经任何手续就被征用。因此我也就无从对这一举动表示同意,更谈不上准许任何人在我不在时进入我的办公室、触碰我的东西。当天我就给UER主任Lefranc先生(Lefranc)打了电话,把情况告诉他;此事(看来)出自Charles夫人的主动,这通电话似乎也证实了这一点。我向Lefranc先生言明:这种做法令我震惊;以如此粗暴的形式调换办公室,要我同意,绝无可能;我期待我的东西在最短期限内物归原位。他向我保证,会去办妥该办的事。

也就在同一天——5月21日,星期二——我写信给Charles夫人,告诉她:我把我办公室这场来得不是时候的“清空”视作滥用权力,并且把它感受为一种暴力;我期待她作出原原本本的解释,期待她毫无保留地道歉。否则,我将把问题提交大学理事会,由它裁定:对USTL的一位教师使出这类手段,是否应当被当作理所当然之事。

今天来到USTL,我得以证实:Charles夫人认为无须回复我的信(信的副本我也送交了Cano与Lefranc两位先生)。Lefranc先生同样认为无须就下述事实向我作任何解释:自从他向我保证会办妥该办的事、让我的东西回到办公室以来,一个星期已经过去,我的办公室依然空空如也。那些被席卷而去的东西如今安在,他也好,Charles夫人也好,都觉得无须告诉我。我辗转从几位秘书那里得知:这些东西存放在她们之中某一位的办公室里。此外,我曾在会议室里与Charles夫人打了个照面,她向我担保,她不过是执行UER主任Lefranc先生的指示,并请我就此事去找他,说这事与她不相干。在事态了结之前,Nguiffo~Boyom先生慨然愿意与我合用他的办公室。

% ---------------- p.831 ----------------
\pmark\origpage{1096}也许,只有我一个人觉得这里面有什么不对劲——一种暴力,一种蔑视;诚然,被人摆出这般架势、不经任何手续扫地出门的,也的确只有我一个。(倘若除我之外还有谁认为,这可不是他愿意在USTL工作的那种气氛,那就请他站出来让我知道,我会由衷高兴……\fnote{834}{(*)}{不言而喻,这样一个举动,对我来说只有一种情况下才有意义:大家明白,它使签署者本人担责,并准许我公开提及此事。})。至于我,我认为,在这场“误会”(姑且借用我一位同事那个可爱的委婉说法)之后,召开一次UER会议,让主任Lefranc先生和Charles夫人有机会就他们的意图与动机作出交代,也让UER的教师们说说,他们是否认为这类做法属正常(当它们落到别人头上的时候……)——这算不得什么奢侈。

我在USTL十二年,常常有机会领教Lefranc先生助人时的善意、热忱与高效——对此我心怀感谢。正因如此,眼见他甘当他人手中的工具,纵容一种专断与蔑视的气氛确立起来,我要收回对他的信任,遗憾也就更添一层。从现在起,我请求他要么切实承担UER主任之责,要么辞去职务。我也请求Charles夫人辞去她那UER“房舍负责人”的职务——这套职务,她滥用得津津有味。

敬候您的(或你的)答复

\hfill Alexandre~Grothendieck

又及:我天性乐于助人。去年,应Lefranc先生之请,我曾同意与Lapscher先生(Lapscher)对换办公室;而不久之后有人告诉我,他随后改变了自己的计划。不言而喻,我的同意,并不意味着我准许对我办公室来一场洗劫——无论在当时,还是在任何别的时候。
\end{quote}

% ---------------- p.832 ----------------
\begin{quote}
朗格多克科学技术大学

\hfill 1985年5月30日,星期四

\underline{数学}

\pmark\origpage{1097}J.~Charles夫人 “数学研究所房舍负责人”\\
致A.~Grothendieck先生,数学教授。

\begin{center}亲爱的同事:\end{center}

1. “数学研究所房舍负责人”的“工作”,从何处开始,到何处为止?这位“负责人”受理数学教师们提出的申请:\\
——或是为一位新任的教师(或研究人员)安排办公室\\
——或是为一位已有办公室的教师(或研究人员)另行安排办公室。

后一类申请,一般都举有工作上的理由:把同一研究组的成员归拢到一处。

这位“负责人”于是着手研究各种可能,首先是同UER~5的主任一道研究——数学研究楼的房舍,其正式委派的管理人,正是由USTL校长指定的那位主任。随后他同有关人士一起寻找可行的方案;调整须待各方谈妥之后方才实施。

2. 近年来照此办成的事情:\\
——几何研究组成员的归拢\\
——力学研究组成员的归拢

3. 这项“工作”中遇到的困难:\\
——实际接触到的每一个人,几乎都觉得自己的办公室归自己“所有”\\
——要强迫什么人“更换”办公室,看来绝无可能。

4. 我接到的最近一份申请,以及为所提问题寻求“方案”的经过评估:\\
——Lapscher(LAPSCHER)教授提出的申请:把Lapscher先生本人、他秘书的办公室、Micali先生归并到同一层。\\
——设想的第一种方案:在三楼与四楼之间交换办公室,好让“申请人”们归并到四楼。这次交换尤其涉及Grothendieck先生与Therond(THEROND)先生。Grothendieck先生经UER~5主任亲自接洽,\dmark{}\origpage{1098}向后者言明:\textbf{办公室在哪儿对他全然无所谓,只要有一间就行}。相反,Therond先生一度表示同意,随后却拒绝了任何搬迁。\\
——设想的第二种方案:我随后请Lapscher先生自己去联络他的同事,以提出另一种方案;这一点UER~5主任向他确认过。他的奔走情形,他不时向我们通报:五间办公室的“占用者”都同意作一次轮换;Grothendieck先生的同意,则出自他同UER~5主任的那次谈话。\\
——第二种方案的实施:UER~5主任得知这一致同意之后,为拟议中的办公室调整开了“绿灯”。

% ---------------- p.833 ----------------
Lapscher先生曾向我谈起钥匙的问题——搬迁启动而未完结的那段期间怎么办;我向他指出:\\
——新钥匙八成一把也没有,\\
——这场搬迁本可由全体有关者一齐动手,几个钟头做完,拖长它,在我看来并不可取。

Lapscher先生随后告知我:Grothendieck先生办公室里的物件,已被搬进他未来的办公室;这件事办成之前,未能先联系上Grothendieck先生。

应当指出:Grothendieck先生住处远在蒙彼利埃之外,目前又借调在法国国家科学研究中心(CNRS);事实上,要同他取得联系非常困难。

5. 以“负责人”的身份,我对这件或许可以称作“冲突”的事情的印象:\\
——我曾有机会向Grothendieck先生言明:我以UER~5的名义行事,不能亲自答复他的信;因此他应当向UER~5的主任索求答复。继这封发给大家的第二封信之后,我认为自己必须走出曾自我施加的“克制义务”。\\
——至少在搬动有关人士的东西之前先通报一声,这在我看来本是可取的\\
——把搬迁压缩到至多半个工作日之内,这在我看来本也是可取的。\\
——所设想的方案在我看来是妥当的:它丝毫改变不了有关各位的办公室占用率。

我不等待答复。

先生、亲爱的同事,请接受我最诚挚的问候。

附记:此信副本送下列各方知悉\\
——收讫Grothendieck先生1985年5月28日来信的全体数学教师;\\
——UER~5主任,他另外还收到Grothendieck先生1985年5月21日写给我的那封信的副本;\\
——USTL临时行政负责人,他已收到5月21日那封信的副本,我再随函附上5月28日那封信的副本。
\end{quote}

% ---------------- p.834 ----------------
\begin{quote}
朗格多克科学技术大学\\
数学研究所\\
\pmark\origpage{1099}数学研究所

1985年6月6日(星期四)18时会议纪要。

出席者:Auberson先生、Charles夫人、Ciulli先生、Contou-Carrère先生、Cuer先生、De~Lima先生、Delobel先生、De~Robert先生、Grothendieck先生、Hocquemiller先生、Escamilla先生、Hubert~Coulin小姐、Lefranc(LEFRANC)先生、Loupias先生、Meden夫人、Molino先生、Pierrot夫人、Pinchard先生、Saint~Pierre先生、Voisin小姐

经讨论后,到会者(19人)以16票赞成、3票弃权通过下列文本:

“数学教师们就Grothendieck先生的东西被搬动时所处的不可接受的情形,向Grothendieck先生致歉。他们共同承诺,警惕这类令人遗憾的事件不再发生。尤其必须明确:一间办公室的钥匙,未经占用者明确同意,任何人不得动用。”

\begin{center}
Lefranc先生\\
主任
\end{center}
\end{quote}

% ---------------- p.835 ----------------
\begin{quote}
\pmark\origpage{1100}朗格多克科学技术大学\\
数学研究所\\
一场“误会”的尾声\\
致数学楼旧日共事的同事们(教学人员与技术人员、第三阶段学生)书\\
Alexandre~Grothendieck

\hfill ……1985年6月7日

亲爱的同事:我在这里为我的办公室遭洗劫一事写下尾声——此事已在我5月28日的信里提起。那封信当时只寄给了数学教师,而它同样关乎、并以同一名义关乎在数学楼里占有一间办公室的每一个人,无论男女。我当初没有把信同时寄给技术人员和第三阶段的学生,是出于疏忽,也失之欠考虑——我(仓促地)以为,那样做等于给这桩插曲一个它并不应得的扩展。我在此向有关各位诚挚致歉——更何况,他们中有好几位(照理本该是“不知情者”……)曾向我表示同情,令我感动。大概也是因为这个疏忽,昨天专门讨论这桩插曲的UER全体大会,只限于“UER~5成员”参加。

在许许多多别的事情之外,这桩插曲还让我领教到:这类事在UER~5并非头一回发生——只是头一回,被瞄准的是一位“A级教师”。昨天表决通过的那份虔诚的决议,能否阻止这类事件重演——在普遍的冷漠之中(一如从前),尤其是对非在编教师、对第三阶段学生的冷漠——我不知道。我会留意向Mori夫人与Moure 夫人核实:她们是否确实收到UER主任的指示——今后无论以什么借口,都不得把任何一间办公室的钥匙交给任何人,也不得为任何人动用,除非得到其中一位占用者的明确许可。

我上一封信,结尾用的是这样几个字:“敬候您的(或你的)答复”。作为对这一期待的回应,我收到三份同情与声援的表示。它们来自Louis~Pinchard、Pierre~Molino与Christine~Voisin。同样,我还收到一份同一方向的表示,出自第三阶段学生Philippe~Delobel——他(与Christine~Voisin一样)曾随我读过DEA(深入研修文凭)。昨天有几位第三阶段学生出席全体大会,正是出于他的倡议。对他,也对我方才提到的所有那些人——他们(不含糊、不闪避)向我表示了声援——我乐于在此表达我的敬重与感谢。像这类的“严酷”经历,果实之一便是让人认出自己的朋友——倘若运气好,还有朋友可认……

我还收到另一封回应我的信,出自一位同事,他分明为发生的事情而喜不自胜,还借此机会温文尔雅地拿我开涮。这是我收到的这一方向的唯一回声。其余诸人那里,不少是这一位的全然冷漠、那一位的局促不安(其间不止一次,我触到一种没有说出口的恐惧:怕给人留下坏印象,从而断送晋升的机会,或危及本已朝不保夕的位置)。在那些人当中——凡是激动到肯为此劳神出席那场全体大会的人(大会是最后一刻偷偷摸摸才召集的,虽然一星期前就已排定……)——我感受最深的,是一种处心积虑、打定主意要把鱼搅浑的心计,口口声声“人人都和善,个个都可爱”。末了(在足足三刻钟的空谈之后),大家转而拿那个众望所归的“坏蛋”问罪,那位缺席者(多么凑巧),Lapscher先生——(据人们方才刚刚暗示出来的意思)\dmark{}\origpage{1101}发起这场偷袭的,正是他。至于走到点名道姓追究他那一步,那是谁也没有打算的,可怜的人——当然,对其他任何一个人,也同样如此,这不必说。

% ---------------- p.836 ----------------
在这桩洗劫插曲中以这种那种名分牵连在内的各“负责人”那里:一位Lapscher的无耻蛮横,一位Charles夫人“以取乐为务”的粗鄙(她在既成事实面前把这场偷袭遮掩过去,还添上她自酿的那一份放肆),以及USTL临时行政负责人Cano先生的失礼——我写信告知他事态原委,请他把问题提交大学理事会,他竟免作任何答复——这一切都令我震惊。然而最甚于一切的,是UER~5主任Lefranc先生那模棱两可的态度,令我错愕,令我痛心。从5月20日星期一(那一天,我把刚发现的情况和我对此事的感受告知了他)直到昨天,他始终不认为适宜之举,或者是向我说明究竟发生了什么,或者是毫不含糊地同一位Lapscher的强盗行径、一位Charles夫人的粗鄙划清界限。自始至终,他竭力维持那场倒霉“误会”的虚构,硬是给那些行径——在我看来无法容忍的行径——装出一副无关痛痒、甚至堪可敬重的模样。为的是不叫任何人难受,可以断言,他选择了“又护山羊又护白菜”:(多多地)照顾山羊,(稍稍地)照顾白菜。

在其他种种迹象之列,我也好好记下了一笔:我原以为可以列入朋友之数的许多人(其中有三位曾是我的学生)所守住的沉默;这一位招摇的漠然,那一位的窘迫不安,还有一位甜腻腻的欢欣雀跃。还有一位Micali的沉默(他是这场偷袭的共同受益人;而几年前,他就已有充分的机会确信:招得Charles先生与夫人不快,会带来多少麻烦……),以及Brun小姐的殷勤卖力——她领一位Lapscher的吩咐,客串了一回雇佣锁匠兼搬家工(等到这场行动的性质再无可疑之时,连一句遗憾的话也没有)。

在这一切的底子上——昨天我又见到那间十二年来一直是我的办公室,这一回它成了战场:我的东西(连家具一并)被七手八脚重新堆作一团(距那场“闪电”偷袭,足足过去了十五天……)——如今我已没有心绪再在里面重新安顿下来。人们向我担保,同样的插曲在我身上重演的可能性微乎其微;何况我尽可以先发制人,把那把迄今托付给Mori夫人与Moure 夫人的第二把钥匙收归己有。但只要物质上办得到——尤其在我借调法国国家科学研究中心的整个期间——我宁愿从今往后放弃在USTL使用办公室,不加争斗地把地方让给Lapscher们、Charles们以及各色同好。

如果能够避免,我不会再在USTL执教。在那儿,我自始至终——可以肯定——不过是个异乡人,一个故乡在别处的人:无论我对待数学的路数,还是我对待教学的路数,抑或我的生活方式,皆是如此。大学这个小天地所能教给我的,我想我已经学到了;而最后一“卷”,便是这桩插曲教给我的一切——它刚刚在一片皆大欢喜中落幕。我刚刚参与的这场UER~5会议,很可能是我的最后一场;这封信,也很可能是我有机会写给你们(或写给你)的最后一封。而这一回,我不等待答复。

\hfill Alexandre~Grothendieck
\end{quote}
